\section{$\neexp$ preliminaries}

\subsection{Introspection games}

In this section, we introduce introspection games.
These are games in which, rather than the verifier sampling the questions,
the provers sample them instead.

\begin{definition}[Introspection games]
An \emph{introspection game} is played between two provers Alice and Bob
in which Alice returns two strings~$x_A$ and~$a$
and Bob returns two strings~$x_B$ and~$a'$ (the verifier does not specify a question).
Here, $x_A$ and~$x_B$ are interpreted as Alice and Bob's ``share" of the jointly sampled ``question" $x = (x_A, x_B)$,
and~$a$ and~$a'$ are interpreted as their ``answers".
The verifier then applies its \emph{evaluation function} $V$ to the answers
and accepts if $V(x_A, x_B, a, b)=1$.
\end{definition}

The following three facts show that we can convert between strategies for a ``normal" game and strategies for an introspective version of the game.
This allows us to prove soundness results for the ``normal" game and ``bootstrap" them up to the introspective game as well.

\begin{fact}\label{fact:introspective-outrospective}
This fact concerns two games and two strategies.
\begin{enumerate}
\item Let $\game_{\mathrm{intro}}$ be the introspective game with evaluation function $V$.
Consider a strategy $\calS_{\mathrm{intro}}$ for Alice and Bob with shared state $\ket{\mathrm{intro}} = \ket{\mathrm{question}}\otimes \ket{\mathrm{answer}}$
in which Alice and Bob's measurements are given by
\begin{equation*}
\{\mathsf{P}_{x_A} \otimes A^{x_A}_a\}_{x_A, a}, \quad \{\mathsf{Q}_{x_B} \otimes B^{x_B}_{a'}\}_{x_B, a'},
\end{equation*} 
respectively.
Write $\calD$ for the distribution on outcomes $(x_A, x_B)$
when the measurement $\{\mathsf{P}_{x_A} \otimes \mathsf{Q}_{x_B}\}_{x_A, x_B}$ is performed on $\ket{\mathrm{question}}$.
\item Let $\game$ be the ``normal" game played as follows: sample $\bx = (\bx_A, \bx_B) \sim \calD$. Distribute the questions as follows:
	\begin{itemize}
	\item[$\circ$] Alice: give $\bx_A$; receive~$\ba$.
	\item[$\circ$] Bob: give $\bx_B$; receive~$\bb$.
	\end{itemize}
	Accept if $V(\bx_A, \bx_B, \ba, \bb) = 1$.
	Write $\calS$ for the strategy with shared state $\ket{\mathrm{answer}}$ in which Alice's strategy is $\{A^{x_A}_a\}_a$
	and Bob's strategy is $\{B^{x_B}_{a'}\}_{a'}$.
\end{enumerate}
Then $\valstrat{\game}{\calS} = \valstrat{\game_{\mathrm{intro}}}{\calS_{\mathrm{intro}}}$.
\end{fact}

\begin{fact}\label{fact:introspectivize-it}
Let $\{A^x_a\}_x$ and $\{B^x_a\}_x$ be measurements such that
\begin{equation}\label{eq:outro-approx}
A^x_a \otimes I_{\reg{Bob}} \approx_\delta B^x_a \otimes I_{\reg{Bob}}
\end{equation}
on state $\ket{\mathrm{answer}}$ and distribution~$\calD$. Next, let $\{\mathsf{Q}_x\}_x$ be a measurement and $\ket{\mathrm{question}}$ be a bipartite state
such that the distribution of measurement outcomes produced by measuring $\{\mathsf{Q}_x \otimes I_{\reg{Bob}}\}_x$ on $\ket{\mathrm{question}}$ is $\calD$.
Then
\begin{equation}\label{eq:intro-approx}
(\mathsf{Q}_x \otimes A^x_a)_{\reg{Alice}} \otimes I_{\reg{Bob}} \approx_{\delta} (\mathsf{Q}_x \otimes B^x_a)_{\reg{Alice}} \otimes I_{\reg{Bob}}
\end{equation}
on state $\ket{\mathrm{question}}\otimes
\ket{\mathrm{answer}}$. Moreover, if $\mathsf{Q}_x$ is a projective
measurement, then the reverse implication holds:
if~\Cref{eq:intro-approx} holds on $\ket{\mathrm{question}} \otimes
\ket{\mathrm{answer}}$, then~\Cref{eq:outro-approx} holds on the state $\ket{\mathrm{answer}}$.
\end{fact}

\begin{proof}
First, we show the forward implication. By definition, we want to bound
\begin{align*}
&\sum_{x, a} \Vert (\mathsf{Q}_x \otimes A^x_a \otimes I_{\reg{Bob}} - \mathsf{Q}_x \otimes B^x_a \otimes I_{\reg{Bob}}) \ket{\mathrm{question}}\otimes \ket{\mathrm{answer}}\Vert^2\\
=& \sum_{x, a} \Vert (\mathsf{Q}_x \otimes I)_{\reg{question}} \otimes (A^x_a \otimes I -   B^x_a \otimes I)_{\reg{answer}} \ket{\mathrm{question}}\otimes \ket{\mathrm{answer}}\Vert^2\\
=& \sum_{x, a} \bra{\mathrm{question}}\otimes \bra{\mathrm{answer}}
	(\mathsf{Q}_x \otimes I)^2_{\reg{question}} \otimes (A^x_a \otimes I -   B^x_a \otimes I)_{\reg{answer}}^2 \ket{\mathrm{question}}\otimes \ket{\mathrm{answer}}\\
\leq& \sum_{x, a} \bra{\mathrm{question}}\otimes \bra{\mathrm{answer}}
	(\mathsf{Q}_x \otimes I)_{\reg{question}} \otimes (A^x_a \otimes I -   B^x_a \otimes I)_{\reg{answer}}^2 \ket{\mathrm{question}}\otimes \ket{\mathrm{answer}}\\
=& \E_{\bx} \sum_{a} \bra{\mathrm{answer}}
	(A^{\bx}_a \otimes I -   B^{\bx}_a \otimes I)^2 \ket{\mathrm{answer}}\\
=& \E_{\bx} \sum_{a} \Vert (A^{\bx}_a \otimes I -   B^{\bx}_a \otimes I)^2\ket{\mathrm{answer}}\Vert^2.
\end{align*}
But this is at most $\delta$ by assumption. Now, for the reverse
implication, note that if $\mathsf{Q}_x$ is projective, then the
inequality above becomes an equality.
\end{proof}

\begin{fact}\label{fact:introspective-consistency}
Let $\{A^x_a\}_x$ and $\{B^x_a\}_x$ be measurements such that
\begin{equation}\label{eq:outro-approx}
(A^x_a)_{\reg{Alice}} \otimes I_{\reg{Bob}} \consistency_\delta I_{\reg{Alice}} \otimes (B^x_a)_{\reg{Bob}}
\end{equation}
on state $\ket{\mathrm{answer}}$ and distribution~$\calD$. Next, let $\{\mathsf{Q}_x\}_x$ be a measurement and $\ket{\mathrm{question}}$ be a bipartite state
such that
\begin{equation*}
(\mathsf{Q}_x)_{\reg{Alice}} \otimes I_{\reg{Bob}} \consistency_\delta I_{\reg{Alice}} \otimes (\mathsf{Q}_x)_{\reg{Bob}}
\end{equation*}
on $\ket{\mathrm{question}}$.
Furthermore, suppose that the distribution of measurement outcomes produced by measuring $\{(\mathsf{Q}_x)_{\reg{Alice}} \otimes I_{\reg{Bob}}\}_x$ on $\ket{\mathrm{question}}$ is $\calD$.
Then
\begin{equation}\label{eq:intro-approx}
(\mathsf{Q}_x \otimes A^x_a)_{\reg{Alice}} \otimes I_{\reg{Bob}} \consistency_{\delta} I_{\reg{Alice}} \otimes (\mathsf{Q}_x \otimes B^x_a)_{\reg{Bob}}
\end{equation}
on state $\ket{\mathrm{question}}\otimes
\ket{\mathrm{answer}}$.
\end{fact}

\subsection{Subroutines and superregisters}

In the next few sections, we will design a set of protocols
to be used as a subroutine for our main $\neexp$ protocol.
In doing so, we will encounter the following notational difficulty:
a subroutine~$\game$ might be a $\lambda = (k, n, q)$-register game,
whereas the overall protocol which calls it might be a $\mu = (\ell, m, q)$-register game.
When~$\lambda$ is not equal to~$\mu$, how can we use~$\game$?
We will consider two answers to this question.
In both of them, we will consider the case when all the register field sizes are the same value ``$q$",
as this is the case relevant to our application.

\begin{notation}
First, the registers in~$\lambda$ might appear as a subset of the registers in $\mu$.
In this case, we will specify an injection $\kappa : [k] \rightarrow [\ell]$
such that $n_i = m_{\kappa(i)}$.
Given a string $W = (W_1, \ldots, W_k)$, we write $\kappa(W)$
for the length-$\ell$ string with $W_{\kappa(i)}$ in coordinate~$i$, for each~$i$,
and the ``hide" symbol $H$ in the remaining coordinates.
Similarly, given string $a = (a_1, \ldots, a_\ell)$,
we write $\kappa^{-1}(a)$ for the length-$k$ string with $a_{\kappa(i)}$ in its $i$-th coordinate.
Then \emph{playing $\game$ on registers $\kappa(1), \ldots, \kappa(k)$} means to do the following.
\begin{enumerate}
\item Draw $(\bx, \bx')$ from $\game$.
\item Send $(\kappa(\bx_1), \bx_2)$ to Alice and $(\kappa(\bx_1'), \bx_2')$ to Bob.
\item Receive $\ba, \ba'$. Accept if $\game$ accepts on the answers $(\kappa^{-1}(\ba_1), \ba_2)$ and $(\kappa^{-1}(\ba_1'), \ba_2')$.
\end{enumerate}
\end{notation}

\begin{notation}
Second, the registers in~$\lambda$ might appear as the concatenation of register in $\mu$.
In this case, we will specify concatenation lengths $c(1) + \cdots + c(k) = \ell$ such that
$n_1 = m_1 + \cdots + m_{c(1)}$, $n_2 = m_{c(1) + 1} + \cdots + m_{c(1) + c(2)}$.
Pictorially, the first register in $\lambda$ will be created as the following concatenation:
\begin{equation*}
\underbrace{
\ket{\mathrm{EPR}_q^{n_1}} \otimes \ket{\mathrm{EPR}_q^{n_2}} \otimes \cdots
\otimes \ket{\mathrm{EPR}_q^{n_{c(1)-1}}} \otimes \ket{\mathrm{EPR}_q^{n_{c(1)}}}. 
}_{\ket{\mathrm{EPR}_q^{n_1 + \cdots + n_{c(1)}}}}
\end{equation*}
We refer to these concatenations of registers as \emph{superregisters}.
A Pauli basis query $W \in \{X, Z, \hideq, \noop\}$ to a given superregister~$R$ can be simulated as follows:
\begin{enumerate}
\item  Implement each Pauli basis query $W$ by sending~$W$ to
  each register $r_{i_1}, \ldots, r_{i_{c}}$ in the superregister.
\item If $W \in \{X, Z\}$, the prover measures $\tau^W_{u_i}$ on each
  register $r_i$ in $R$, and the verifier receives the outcomes
$
\bu_{i_1} \in \F_q^{m_{i_1}}, \ldots, \bu_{i_{c}} \in \F_q^{m_{i_{c}}},
$
concatenated as $\bu = (\bu_{i_1}, \ldots, \bu_{i_c})$.
\item If $W = \hideq $, the prover performs $\oktimes{I}{c}$ on
  the registers in the superregister, and verifier receives~$c$ consecutive $\varnothing$'s, interpreted as a single $\varnothing$.
\item If $W = \bot$, the prover may perform any measurement it likes
  on the registers in the superregister.
  The verifier receives~$c$ consecutive $\varnothing$'s, interpreted as a single $\varnothing$.
\end{enumerate}
\end{notation}

The game~$\game$ will usually be proven sound against $\lambda$-register strategies,
but in our cases it will be straightforward to extend this soundness to $\mu$-register strategies
in the case when~$\game$ is applied as a subroutine as detailed above.
For example, suppose we know that a strategy~$A$ which passes~$\game$
with probability $1-\eps$ must satisfy
$(A_a)_{\reg{Alice}} \otimes I_{\reg{Bob}} \approx_{\delta} (B_a)_{\reg{Alice}} \otimes I_{\reg{Bob}}$.
Then it is straightforward to derive that if~$\game$ is played as a subroutine on the second register of~$\mu$
(this is the case when $k = 1$ and $n_1 = m_2$), and if~$A$ passes the subroutine with probability~$1-\eps$, then
\begin{equation*}
(A_a)_{\reg{Alice}} \otimes I_{\reg{Bob}} \approx_{\delta} (I_{\reg{1}} \otimes (B_a)_{\reg{2}} \otimes I_{\reg{3, \ldots, \ell}})_{\reg{Alice}}.\otimes I_{\reg{Bob}}
\end{equation*}
Likewise, suppose $\game$ is played as a subroutine on one superregister consisting of all the registers of~$\mu$
(this is the case when $k = 1$ and $n_1 = m_1 + \cdots + m_\ell$).
If $A$ passes the subroutine with probability~$1-\eps$, then
\begin{equation*}
(A_a)_{\reg{Alice}} \otimes I_{\reg{Bob}} \approx_{\delta} ((B_a)_{1, \ldots,\ell} )_{\reg{Alice}} \otimes I_{\reg{Bob}}.
\end{equation*}
For our applications, it will be straightforward to
extend the soundness of our games to the case when they are played as subroutines,
and we will leave this step implicit in our proofs.

%%% Local Variables:
%%% mode: latex
%%% TeX-master: "main"
%%% End:
%%%---------- close: big-low-degree-prelims

%%%%%%%%%%% jump to big-low-degree
%%%---------- open: big-low-degree
%!TEX root = main.tex

